\chapter*{Zur Wiederkehr von \texorpdfstring{$2/3$}{2/3}: Eine Antwort auf Austins algebraische Notiz}
\addcontentsline{toc}{chapter}{Antwort auf Austins algebraische Notiz zu $2/3$}

\begin{flushright}\small
Johann Pascher \\
in Antwort auf Peter Austin, \emph{A Purely Algebraic Short Note on the
Repeated Appearance of $2/3$} (Mai 2026) \\
für die IPI-Korrespondenzgruppe --- Mai 2026
\end{flushright}

\section{Zweck und Haltung}

Peter, danke für die Notiz. Sie ist sorgfältig, die Algebra ist korrekt,
und sie rekonstruiert den FFGFT-Leptonensektor zutreffend aus den
öffentlichen Zusammenfassungen. Was folgt, ist weder eine Verteidigung noch
eine Bitte um Bestätigung: es sind die internen Ableitungen, nach denen die
Notiz im Grunde fragt. Wo deine Rekonstruktion bei den öffentlichen Zahlen
endet, trägt FFGFT die Struktur bereits eine Ebene tiefer, und an zwei
Stellen gelangt deine Notiz unabhängig zu Schlüssen, die das Modell aus
sich heraus erreicht hat. Ich verwende durchgehend deine Abschnittsnummern.

Zur Klarstellung: die kanonischen geladenen Leptonen-Koeffizienten sind
\[
  (r_e, r_\mu, r_\tau) = \left(\tfrac{4}{3},\, \tfrac{16}{5},\, \tfrac{25}{9}\right),
  \qquad
  (p_e, p_\mu, p_\tau) = \left(\tfrac{3}{2},\, 1,\, \tfrac{2}{3}\right),
\]
mit $m_i = r_i\,\xipar^{p_i}\,v$ und $\xipar = 4/30000$. Deine Notiz
verwendet genau diese Werte; sie sind aktuell.

\section{Der $\epsilon$-Deformations-Diskriminator (dein \S6)}

Ich akzeptiere den Diskriminator ohne Vorbehalt. Übereinstimmung bei $D=3$
ist keine Identität des Mechanismus, und die Fortsetzung $D = 3+\epsilon$ ist
das richtige Instrument, um die Wege zu trennen. Das ist ausdrücklich
festzuhalten, weil FFGFT \emph{dieselbe} Trennung von innen erreicht, auf
einem anderen Weg (siehe \S5): der Leiter-Wert der Koide-Größe und die
Gleichgewichtsbedingung $A=B$ fallen nicht exakt zusammen --- sie treffen
sich \emph{nahe} $2/3$ und trennen sich, genau wie es deine
$\epsilon$-Analyse vorhersagt. Wir sind uns über die Methode nicht uneinig;
wir sind uns über das Ergebnis einig.

\section{Die Leiter ist geometrisch, nicht arithmetisch (dein \S11)}

Dein \S11 hat recht, dass $\left(\tfrac{3}{2}, 1, \tfrac{2}{3}\right)$ keine
gleichförmige $1/3$-Leiter ist: die Sprünge betragen
\[
  p_e - p_\mu = \tfrac{1}{2}, \qquad p_\mu - p_\tau = \tfrac{1}{3}.
\]
Die Auflösung ist, dass die Leiter überhaupt nicht arithmetisch ist. Sie ist
eine \emph{geometrische} Folge mit Verhältnis $q = 2/3$:
\[
  p_{n+1} = \tfrac{2}{3}\,p_n,
  \qquad
  \tfrac{3}{2} \xrightarrow{\times 2/3} 1 \xrightarrow{\times 2/3} \tfrac{2}{3}.
\]
Die ungleichen Differenzen $1/2$ und $1/3$ sind dann \emph{Konsequenzen}
eines konstanten Verhältnisses, nicht eigenständige Regeln. Das Elektron
braucht daher keine Ad-hoc-Sonderregel; es ist die Basis-Sprosse der Folge.
Zwei gleichwertige Lesarten desselben Sachverhalts:

\begin{itemize}
  \item \textbf{Verhältnis-Lesart.} $q = 2/3$ ist die Inverse des
        Tetraeder-Geometriefaktors $3/2$ (vier Ecken über drei Kanten pro
        Ecke). Dieselbe Zahl, die die Leiter ordnet, ist numerisch der
        Koide-Wert --- genau deshalb erscheinen beide verknüpft.
  \item \textbf{Komplement-Lesart (dein Weg).} An der $\mu\!\to\!\tau$-Sprosse
        ist der Schritt $1 - 1/D = 2/3$ mit $D = 3$, also
        $p_\tau = 1 - 1/D$. Das verortet FFGFT auf deiner ``complement of
        step''-Straße, wie du korrekt erkannt hast, und \emph{nicht} auf der
        Gleichgewichts-Straße $2/D$.
\end{itemize}

Im harmonischen Bild sitzen die Sprossen auf einem logarithmischen Torus in
Schritten einer musikalischen Terz, $\Delta p = 1/3$. Das Elektron bei
$p_e = 3/2 = 1 + 1/(D-1)$ ist eine \emph{halbzahlige} Sprosse --- eine
geometrische Mitte zwischen Windungsebenen (der Spin-Windungsbeitrag), kein
voller Windungsschritt. Das ist der strukturelle Ursprung des
$1/2$-Sprungs.

\section{Topologisches $D=3$ gegen die fraktale Dimension (dein \S11--\S12)}

Dein \S12 fragt, ob der $1/3$-Schritt an das ganzzahlige $D=3$ oder an
$D_f = 3-\xi$ gebunden ist, wobei Letzteres eine Verschiebung nach unten,
$2/3 - \xi/9$, ergäbe. Die Antwort des Modells ist konkret, und sie löst die
Sorge gerade für den Koide-Fall auf.

Erstens gibt es \emph{zwei} abgeleitete fraktale Dimensionen, nicht eine:
\[
  D_f^{\,\text{Raum}} = 3 - \xi \approx 2{,}9998667
  \quad\text{(lokale Metrik; GRB-Laufzeit)},
\]
\[
  D_f^{\,\text{eff}} \approx 2{,}973
  \quad\text{(kumulativer Skalenfluss, Planck bis elektroschwach)}.
\]
Beide sind \emph{abgeleitet}, nicht angenommen. Der Exponentenschritt ist an
die \emph{ganzzahlige, topologische} Torus-Dimension $D=3$ gebunden; die
Metrik $3-\xi$ tritt anderswo ein (Dispersion), und das effektive $2{,}973$
geht nur in die absolute Normierung über
$K_{\text{frak}} = (D_f^{\,\text{eff}}/3)^{D_f^{\,\text{eff}}/2} = 1 - 100\,\xi$
ein.

\begin{keyresult}[Warum keine $\xi/9$-Korrektur Koide erreicht]
Die Koide-Größe $Q$ ist ein reines Massen\emph{verhältnis}. Jeder
skalenabhängige Faktor --- der Vakuumwert $v$, die fraktale Korrektur
$K_{\text{frak}}$ und jede Verschiebung auf $D_f$-Ebene --- ist allen drei
Leptonen gemeinsam und kürzt sich in $Q$ heraus. Fraktale Korrekturen haften
an \emph{absoluten} Skalen, nie an Verhältnissen. Deine Verschiebung
$2/3 - \xi/9$ aus \S12 wäre ein realer Effekt auf einer absoluten Vorhersage;
sie geht schlicht nicht in $Q$ ein, weil $Q$ die absolute Skala nie sieht.
\end{keyresult}

Für Koide ist die Wahl damit zugunsten deiner Option (a) entschieden: der
Schritt liegt auf dem topologischen $D=3$, und $D_f$ stört das Verhältnis
nicht.

\section{Erzwungen oder nur gefittet? Eine ehrliche Antwort (dein \S10, \S13)}

Das ist die zentrale Frage deiner Notiz, und die interne Antwort von FFGFT
liegt bereits vor --- und sie ist interessanter als ein glattes
``erzwungen''.

\subsection*{Was strukturell erzwungen ist}

Die Massen\emph{verhältnisse} sind nicht gefittet. Zwei Punkte:
\begin{itemize}
  \item $m_\mu/m_e$ folgt auf zwei unabhängigen Wegen, die übereinstimmen:
        geometrisch aus der Leiter, $\frac{12}{5}\,\xi^{-1/2} \approx 207{,}8$,
        und fraktal aus $(D_f^{\,\text{eff}}/3)^{D_f^{\,\text{eff}}/2}\!\cdot f(\xi)
        \approx 206{,}8$. Die Übereinstimmung legt $K_{\text{frak}}$ eindeutig
        fest --- ein Konsistenztest, kein Fit.
  \item In der kompakten Darstellung sind die Koeffizienten keine freien
        rationalen Zahlen: $\tfrac{2}{3} = 3\sqrt{3}/(2\pi\alpha^{1/2})$ und
        $\tfrac{8}{5} = 9/(4\pi\alpha)$, so gebaut, dass $\alpha$ sich exakt
        herauskürzt. Die Rationalzahlen sind $\alpha$-geometrische
        Invarianten.
\end{itemize}

\subsection*{Was die Leiter \emph{nicht} erzwingt}

Sie erzwingt die Gleichgewichtsbedingung $A=B$ nicht exakt. Setzt man die
FFGFT-Massen in die Koide-Formel ein, erhält man
\[
  Q_{\text{FFGFT}} \approx 0{,}6677,
  \qquad
  \Delta_K = 3Q_{\text{FFGFT}} - 2 \approx +0{,}003,
\]
d.\,h.\ deine Residuum-Diagnose $\Delta_K = B/A - 1$ ist klein, aber von
null verschieden: $A = B$ gilt auf etwa $0{,}16\,\%$ in $Q$, nicht exakt.
Das ist beabsichtigt und ehrlich. Der Grund ist strukturell und ist dein
eigener Punkt in anderer Gestalt: der Koide-Nenner
$\big(\sum \sqrt{m_i}\big)^2$ erzeugt Kreuzterme mit Exponenten
\[
  \xi^{5/4}, \quad \xi^{13/12}, \quad \xi^{5/6},
\]
von denen keiner ein Vielfaches von $1/3$ ist. Sie liegen \emph{außerhalb}
des FFGFT-Strukturraums $\{p\} \subset \tfrac{1}{3}\mathbb{Z}$; sie sind
Artefakte der Quadratwurzel-Komposition, keine Torus-Moden. Exaktes
$Q = 2/3$ verlangte
\[
  m_e + m_\mu + m_\tau = 4\big(\sqrt{m_e m_\mu} + \sqrt{m_e m_\tau} + \sqrt{m_\mu m_\tau}\big),
\]
was nicht bei $\xipar = 4/30000$ gilt, sondern bei $\xi_{\text{Koide}} \approx
1{,}372\times 10^{-4}$, etwa $2{,}9\,\%$ größer. Mit anderen Worten: die
Gleichgewichtsbedingung ist eine algebraische Nebenbedingung, die
\emph{unabhängig} von der Torus-Geometrie ist --- genau das, was deine
$\epsilon$-Deformation sagt: die Leiter-Straße und die Gleichgewichts-Straße
sind verschiedene Mechanismen, die sich nahe $2/3$ treffen und trennen.

\begin{keypoint}[Die ehrliche Position]
Die Leiter \emph{sagt} $Q$ auf $0{,}16\,\%$ aus $(r_i, p_i, \xi)$
\emph{voraus}, ohne Koide-Eingabe. Der exakte experimentelle Wert $Q = 2/3$
\emph{bestätigt} dann die Exponentenstruktur aus einer unabhängigen
Richtung; er ist weder in FFGFT hineingefittet noch algebraisch von ihr
erzwungen. Das kleine Residuum ist real und liegt genau dort, wo deine
Kreuzterm-Analyse es vorhersagt.
\end{keypoint}

\section{Drei Straßen und die Abbildung (dein \S13)}

Ich stimme deiner Rahmung zu: die Koide-Projektorform, die
FFGFT-$\xi$-Leiter und Pauls $24$-Zellen-Rangaufspaltung sind drei
verschiedene Primitivobjekte, und eine Beziehung zwischen ihnen muss durch
eine Abbildung gezeigt werden, nicht durch eine gemeinsame Zahl. Die
FFGFT-zu-Koide-Abbildung ist explizit,
\[
  (r_i, p_i, \xi)\ \longmapsto\ x_i = \sqrt{r_i}\,\xi^{p_i/2}\,\sqrt{v}
  \ \longmapsto\ (A, B)\ \longmapsto\ Q,
\]
und dein trigonometriefreies Ziel ist in unserer Notation identisch mit der
erzeugten Größe
\[
  Q_{\text{FFGFT}} =
  \frac{\sum_i r_i\,\xi^{p_i}}{\big(\sum_i \sqrt{r_i}\,\xi^{p_i/2}\big)^2},
\]
da sich $\sqrt{v}$ herauskürzt. Wir berechnen dasselbe Objekt. Der einzige
mögliche Streitpunkt ist nicht die Abbildung, sondern die Behauptung der
Exaktheit: die Abbildung schickt die Leiter auf $Q \approx 0{,}6677$, nicht
glatt auf $2/3$, aus dem Grund in \S5.

\section{Schluss}

Kurz also: deine Methode ist richtig; deine Verortung von FFGFT auf der
Komplement-Straße ist richtig; die Leiter-Asymmetrie löst sich als
geometrische Folge mit Verhältnis $2/3$ auf; die $D_f$-Verschiebung erreicht
Koide nicht, weil $Q$ ein Verhältnis ist; und die Frage ``erzwungen oder
gefittet'' hat die ehrliche Antwort, dass die Leiter $Q$ auf $0{,}16\,\%$
vorhersagt und der exakte Wert eine Bestätigung auf einer unabhängigen
Straße ist --- was deine eigene $\epsilon$-Analyse vorwegnimmt. Falls es
nützlich ist, darfst du dies gern neben deine Notiz stellen; ich denke, die
beiden ergeben gemeinsam die klarere Aussage.

\begin{flushright}\small
--- J.P.
\end{flushright}
